% !TEX TS-program = lualatex
% !TEX encoding = UTF-8 Unicode
\documentclass[ngerman,11pt]{article}
\usepackage{fontspec}
\usepackage[ngerman]{babel}
\usepackage[ngerman]{datetime2}
\usepackage[a4paper, total={16cm, 25cm}]{geometry}
\usepackage{amsmath}
\usepackage{amsfonts}
\usepackage{amssymb}
\usepackage{multicol}
\usepackage{tcolorbox}
\usepackage{enumitem}
\usepackage{url}
\usepackage[colorlinks=true, linkcolor=blue, urlcolor=blue]{hyperref}
\usepackage{listings}
\lstset{
    language=Python,
    basicstyle=\ttfamily\small,
    keywordstyle=\color{black},
    stringstyle=\color{black},
    commentstyle=\color{black},
    breaklines=true,
    frame=single,
    framesep=1mm,
    framextopmargin=1mm,
    framexbottommargin=1mm,
    showspaces=false,
    showstringspaces=false,
    numbers=left,
    numberstyle=\tiny\color{gray}
}
\date{\DTMgermanmonthname{\month} \the\year}

\begin{document}
\hrule
\begin{center}
{\large\bf Python: Programmierkonstrukte - Übungsblatt 1}\\[-0mm]
Studienkolleg der TU Berlin\hfill Till Zoppke\\
Informatik \hfill \DTMgermanmonthname{\month} \the\year\\[2mm]
\end{center}
\hrule
\par\bigskip

\section*{Aufgabe 1: Code-Analyse}

Erklären Sie, was das folgende Programm macht. Beschreiben Sie dabei:
\begin{itemize}[leftmargin=*]
\item Was wird eingegeben?
\item Welche Berechnungen werden durchgeführt?
\item Was wird ausgegeben?
\item Geben Sie ein konkretes Beispiel mit Eingabe und Ausgabe an.
\end{itemize}

\begin{lstlisting}
name = input("Wie heißen Sie? ")
alter = int(input("Wie alt sind Sie? "))

if alter >= 18:
    status = "volljährig"
    jahre_bis_rente = 67 - alter
else:
    status = "minderjährig"
    jahre_bis_volljaehrigkeit = 18 - alter

print(f"Hallo {name}!")
print(f"Sie sind {status}.")

if alter >= 18:
    if jahre_bis_rente > 0:
        print(f"Bis zur Rente sind es noch {jahre_bis_rente} Jahre.")
    else:
        print("Sie sind bereits im Rentenalter oder erreichen es bald.")
else:
    print(f"In {jahre_bis_volljaehrigkeit} Jahren werden Sie volljährig.")
\end{lstlisting}

\section*{Aufgabe 2: Programmlogik}
Was gibt dieses Programm aus, wenn x = 5 eingegeben wird? Erklären Sie das Verhalten:

\begin{lstlisting}
x = int(input("Geben Sie eine Zahl ein: "))
if x > 0:
    print("Positiv")
if x > 10:
    print("Groß")
else:
    print("Klein oder gleich 10")
\end{lstlisting}

\section*{Aufgabe 3: Fehlersuche}

Der folgende Code soll das Alter einer Person berechnen. Finden Sie alle Fehler und korrigieren Sie diese:

\begin{lstlisting}
geburtsjahr = input("In welchem Jahr wurden Sie geboren? )
aktuelles_jahr = 2024
alter = geburts_jahr - aktuelles_jahr
print(f"Sie sind alter Jahre alt.")
\end{lstlisting}

\section*{Aufgabe 4: Quotienten berechnen}
Vervollständigen Sie das Programm, das zwei Zahlen von der Tastatur einliest und ihren Quotienten ausgibt. Ersetzen Sie dazu die \ldots~durch geeigneten Programmcode.

\begin{lstlisting}
# Eingabe
zahl1 = ...
zahl2 = ...

# Fehlerbehandlung bei Division durch 0
if ...
	print ...
else:
	# Berechnung und Ausgabe
	quotient = ...
	print(f"Quotient: ...
\end{lstlisting}

\section*{Aufgabe 5: Schaltjahr prüfen}
Schreiben Sie ein Programm, das prüft, ob ein Jahr ein Schaltjahr ist. Ein Jahr ist genau dann ein Schaltjahr, wenn es durch 4 teilbar ist. Es darf dabei nicht durch 100 teilbar sein. Eine Ausnahme gilt für Jahre, die durch 400 teilbar sind: Diese sind trotzdem Schaltjahre.

Sie können ihr Programm mit diesen Eingaben testen: 1900 (kein Schaltjahr), 2000 (Schaltjahr), 2023 (kein Schaltjahr) und 2024 (Schaltjahr).


\end{document}
